\chapter{Conclusion}
The first set of research questions focus on the methods commonly used for computing orbital ascent trajectories. 

\textbf{How are orbital ascent trajectories currently computed?}
\begin{itemize}
  \item What are the challenges faced when computing an orbital ascent trajectory?
  \item What methods are currently used or being researched to solve this problem?
  \item What are the limitations of methods?
\end{itemize}

The most common methods used to compute orbital ascent trajectories are based on trajectory optimization, a set of mathematical tools used to compute open-loop solutions to optimal control problems. The reason is because these methods are well suited for solving problems that are non-linear, high dimensional and sensitive to errors, a set of traits which are unfortunately true when computing orbital ascent trajectories.

Trajectory optimization methods can be split in two categories, direct and indirect methods. With direct methods the infinite-dimensional trajectory optimization problem is first transcribed to a finite-dimensional approximation which can then be solved using a parameter optimization method. In contrast with indirect methods a surrogate problem is created by analytically deriving the necessary conditions for optimal control and then the new problem is transcribed into a finite-dimensional approximation, which can then be solved using a parameter optimization method.

Indirect methods are more commonly used for computing orbital ascent trajectories, because these tend to be more accurate and provide a better measure of accuracy. However these methods are more complex then direct methods since an analytical model is required and the surrogate problem needs to be derived analytically.

Furthermore during the derivation of the surrogate problem a new set of adjoint variables are added, which also need to be initialized and optimized together with the controls and states, and since these adjoint variables don't have a physical representation, it may be difficult to find appropriate initial values when solving the problem.

To ease the complexity of initializing indirect problems, Pontani \emph{et al.} \cite{PONTANI2014852, pontani2013ascent, pontani2014simple, pontani2014optimal} proposes the use of \acrfull{pso} to numerically determine the control, state and adjoint variables after analytically deriving the surrogate problem. Due to its heuristic nature \acrshort{pso} is able to try a wider range of initial conditions and thus significantly lessens the initialization issues while retaining the accuracy that comes with indirect methods.

The main limitation of trajectory optimization methods is that only open-loop solutions can be obtained. These are solutions that consist of an initial condition and a sequence of controls that need to be applied through time. Some drawbacks of this type of solution are a high sensitivity to initial conditions and lack of guarantee that the solutions converge to the global optimal. Furthermore a stabilizing controller is required to keep the system on the predefined trajectory when applied to real systems.

In contrast closed-loop solutions provide the optimal controls for the entire state space, thus closed-loop solutions are able to guarantee the global optimal solution and can be applied directly to real systems. The drawback is that these solutions are difficult to compute, specially when applied to high dimensional problems.

Originally developed in the 1950's, \acrfull{dp} is a class of methods used to compute closed-loop solutions for optimal control problems. However since these methods require repeated scans of the entire state space, the complexity and computational requirements increases exponentially as more dimensional and states are added to the problem.

These problems can be mitigated or solved by combining \acrshort{dp} with various machine learning techniques to produce what is now known as \acrfull{rl}. Thanks to these techniques, \acrshort{rl} methods are able to learn closed-loop solutions with only partial views of the state space by generalizing between visited states and actions.

The next research questions focus \acrshort{rl} and how it can be used to solve the aforementioned problems. 

\textbf{What is \acrlong{rl}?}
\begin{enumerate}
  \item How does \acrshort{rl} work?
  \item What is \acrshort{rl} capable of?
  \item How has \acrshort{rl} been used so far?
  \item What state of the art technological developments make use of \acrshort{rl}?
\end{enumerate}

\textbf{How can \acrlong{rl} methods be used to solve the orbital ascent trajectory opti­mization problem?}
\begin{enumerate}
  \item What advantages do \acrshort{rl} methods have over existing methods?
  \item What properties of \acrshort{rl} can be used to circumvent the limitations of current orbital ascent trajectory optimization methods?
  \item Are \acrshort{rl} based orbital ascent trajectory computation methods capable of achieving the accuracy of existing methods?
  \item Is an \acrshort{rl} based closed-loop solution able to outperform a precomputed open-loop solution in case of unknown disturbances or errors in the training models.
\end{enumerate}

Both \acrshort{dp} and \acrshort{rl} are used to solve a class of stochastic optimal control problem known as \acrfull{mdp}. These problems consist of two main elements, an agent that makes decisions based on a policy, interacts with its environment and learns. And the environment which consists of everything else. The agent will perform an action which changes the environment and then the environment responds by giving the agent a reward, which it then uses to learn. The goal of the agent is to maximize the cumulative reward it receives through time.

The reward only gives an indication of the immediate effects of the agent's actions, it shows which states are good and bad. These rewards are used to compute the value function. A function that gives an indication of which state-actions pairs are more likely to yield higher rewards in the long run.

Some \acrshort{rl} methods take decisions directly from the estimated value functions, however a different class of methods that estimate the policy directly are better able to handle high dimensional problems with continues action spaces. These are known as policy gradient methods. These methods can still take advantage of an estimated value function to speed up the learning process, these methods are known as actor-critic methods, where the policy estimate is the decision taking actor and the value function is the critic.

Various learning techniques exist for \acrshort{rl}, the main one being used in this thesis is \acrfull{ppo}. The main scheme used in this technique is limiting the difference between two policies during an update within a given trust region. This minimizes the risk of a bad update and allows multiple updates to the policy to be performed using the same set of samples, thus improving sample efficiency \cite{schulman2017proximal, schulman2015trust}.

The use of \acrshort{rl} in spacecraft navigation is relatively new and limited to research settings. The main characteristics of \acrshort{rl} that of interest are the ability to generalize and autonomously adapt to unknown situations. For example Gaudet and Furfaro \cite{gaudet2012robust} successfully explored the use of RL as a control method for hovering in the non-uniform, rotating gravitational field commonly associated with asteroids. 

Hovell and Ulrich \cite{hovell2020deep, hovell2021deep} and Broida and Linares \cite{broida2019spacecraft} have been developing RL based control system for proximity operation guidance that could be applied in situations where the chaser spacecraft may be required to manoeuvre in the proximity of a potentially uncooperative target object.

Gaudet et al. \cite{gaudet2018deep} also demonstrated the use of an \acrshort{rl} agent for planetary powered landing, with the aim to develop a landing system capable to navigate within meters from its target. And finally \acrshort{rl} has also been successfully demonstrated to be a suitable tool to develop efficient control systems for efficient orbital change manoeuvre with low-thrust spacecraft.

The main project for this thesis is the use of \acrshort{rl} to compute closed-loop solutions for orbital ascent trajectories. Although the overall training process is likely more computationally intensive than computing a single solution using trajectory optimization, an already trained agent would be able to compute the optimal or close to optimal solutions online to compensate for disturbances. Thus the main research objective for this thesis is:

\emph{“To obtain a closed-loop orbital accent guidance law capable of generalizing beyond explored environments and compensate for unknown disturbances using Reinforcement Learning.”}

So far most research questions have been answered in this preliminary thesis. The remaining questions that still need to be investigated are:

\textbf{How can \acrlong{rl} methods be used to solve the orbital ascent trajectory opti­mization problem?}
\begin{itemize}
  \item Are \acrshort{rl} based orbital ascent trajectory computation methods capable of achieving the accuracy of existing methods?
  \item Is an \acrshort{rl} based closed-loop solution able to outperform a precomputed open-loop solution in case of unknown disturbances or errors in the training models.
\end{itemize}

To answer these questions the following plan is proposed. The project will consists out of two main programs, a \acrfull{to} based program using the method proposed by Pontani \emph{et al.} \cite{PONTANI2014852, pontani2013ascent, pontani2014simple, pontani2014optimal} and program using a novel \acrshort{rl} based method.

Since the achieved performance is not known until the terminal state, it is necessary to use a method of supervised or imitation learning to help initialize the models used by the \acrshort{rl} agent. The \acrfull{to} based program will be used as the supervisor, thus the first step of the project is to develop it. 

To assess the accuracy of the \acrshort{rl} based program, the results of it can be compared to the results yielded by the \acrshort{to} based program. If accurate both, should yield trajectories with the same or similar performance, with the performance in this case being measured by propellant utilization and accuracy of the orbital insertion manoeuvre end state.

To test the resistance of the \acrshort{rl} based solutions to disturbances, the agents will be trained and tested in stochastic scenarios introducing different types of disturbances. The ability to generalize should help the agents to adapt to unseen environment and thus help compensate for modeling errors or other unknowns. This can be tested by judging the agent's ability to succeed in unknown environments and compare the resulting performance to optimal solutions.
